\section{Introduction}
\label{sec:introduction}

How language and communication emerge among intelligent agents has long been a topic of intense debate. Among the many unresolved questions are: Why does language use discrete structures? What role does the environment play? What is innate and what is learned? And so on. Some of the debates on these questions have been so fiery that in 1866 the French Academy of Sciences banned publications about the origin of human language.

The rapid progress in recent years of machine learning, and deep learning in particular, opens the door to a new perspective on this debate.  How can agents use machine learning to automatically discover the communication protocols they need to coordinate their behaviour?  What, if anything, can deep learning offer to such agents?  What insights can we glean from the success or failure of agents that learn to communicate?

In this paper, we take the first steps towards answering these questions.  Our approach is programmatic: first, we propose a set of multi-agent benchmark tasks that require communication; then, we formulate several learning algorithms for these tasks; finally, we analyse how these algorithms learn, or fail to learn, communication protocols for the agents.

The tasks that we consider are fully cooperative, partially observable, sequential multi-agent decision making problems.  All the agents share the goal of maximising the same discounted sum of rewards.  While no agent can observe the underlying Markov state, each agent receives a private observation correlated with that state.  In addition to taking actions that affect the environment, each agent can also communicate with its fellow agents via a discrete limited-bandwidth channel.  Due to the partial observability and limited channel capacity, the agents must discover a communication protocol that enables them to coordinate their behaviour and solve the task.

We focus on settings with \emph{centralised learning} but \emph{decentralised execution}.  In other words, communication between agents is not restricted during learning, which is performed by a centralised algorithm; however, during execution of the learned policies, the agents can communicate only via the limited-bandwidth channel.  While not all real-world problems can be solved in this way, a great many can, e.g., when training a group of robots on a simulator.  Centralised planning and decentralised execution is also a standard paradigm for multi-agent planning \cite{kraemer2016multi}. For completeness, we also provide decentralised learning baselines. 

To address these tasks, we formulate two approaches.  The first, named \emph{reinforced inter-agent  learning} (RIAL), uses deep $Q$-learning \cite{Mnih:2015} with a recurrent network to address partial observability. In one variant of this approach, which we refer to as  \emph{independent Q-learning}, the agents each learn their own network parameters, treating the other agents as part of the environment. Another variant trains a single network whose parameters are shared among all agents. Execution remains decentralised, at which point they receive different observations leading to different behaviour.

The second approach, which we call \emph{differentiable inter-agent  learning} (\dic), is based on the insight that centralised learning affords more opportunities to improve learning than just parameter sharing.  In particular, while RIAL is end-to-end trainable \emph{within} an agent, it is not end-to-end trainable \emph{across} agents, i.e., no gradients are passed between agents.  The second approach allows real-valued messages to pass between agents during centralised learning, thereby treating communication actions as bottleneck connections between agents.  As a result, gradients can be pushed through the communication channel, yielding a system that is end-to-end trainable even across agents.  During decentralised execution, real-valued messages are discretised and mapped to the discrete set of communication actions allowed by the task.  Because \dic passes gradients from agent to agent, it is an inherently deep learning approach.

Our empirical study shows that these methods can solve our benchmark tasks, often discovering elegant communication protocols along the way. To our knowledge, this is the first time that either differentiable communication or reinforcement learning (RL) with deep neural networks 
have succeeded in learning communication protocols in complex environments involving sequences and raw input images\footnote{This paper extends our own work in~\cite{foerster2016learning}, which proposes a variation of DQN for protocol learning.}. The results also show that deep learning, by better exploiting the opportunities of centralised learning, is a uniquely powerful tool for learning communication protocols. Finally, this study advances several engineering innovations, outlined in the experimental section, that are essential for learning communication protocols in our proposed benchmarks.

